% P10-Evaluationsfeinschliff, 16.09.2026: Aussagegrenzen und Anforderungsbilanz; keine neuen Messwerte.
% B-59, 15.09.2026: datierte Ledger-Nachauswertung, Stützung und Reichweite.
% Beleg: thesis-evidence/w2/ledger-repetitions-20260915/ (Originale unverändert).
% B-58, 15.09.2026: Referenzzahlen nach Fall benannt; Werte unverändert.
% P10-4c, 14.09.2026: Ledger-Herleitung und Nachweisanschlüsse präzisiert.
% M19 (12.09.2026): Leserführung und präzise Ergebnisdarstellung;
% Messwerte und historische Nachweisgrenzen bleiben erhalten.
% M02 (10.09.2026), B-14: Inputverhältnis F/Ledger rund 78 Prozent.
% HERKUNFT: thesis-evidence/w2/token-usage-correction-20260910.json.
% ============================================================
% §7.3 Ergebnisqualität der Ledger-Kette — v02 (Claude, 08.09.)
% NACHREVIEW-STAND: konsolidierte Matchurteile der Adjudikations-
% runde 08.09. (8 Änderungen: Ledger 025/062/086/105/107,
% F 025/089/107; Auslegung G-IE-105 = none beidseitig).
% Struktur: 7.3.1 Ergebnisqualität · 7.3.2 Lücken/Hinweise (a/b/c)
% · 7.3.3 Stufenanalyse · 7.3.4 F-Einordnung. Nenner 109 / Gold v2.
% Leitplanken (Auftrag 08.09.): Hinweis ≠ Erkennung ≠ Behebung ·
% Szenario = bedingte Rechnung · Claim-Zahl ≠ Qualität · keine
% Wichtigkeitsskala aus Lückentypen · kein Gesamtsieger.
% ============================================================
% HERKUNFT: Konsolidierter Auswertungsstand + alle abgeleiteten
% Zahlen = runs/w2-nachreview-konsolidierung.json (Beziehungstabelle
% je Gold-ID; Kopie thesis-evidence/w2/eval/); Entscheidungsrunde =
% kapitel-7/ADJUDIKATION-Auftrag-Kapitel-7-Nachreview.md (beauftragte
% KI-Adjudikation 08.09., getrennt von der 8+3-Historie); Reviews =
% kapitel-7/Review-Coverage-109-Aussagen-8189ea.md + Review-F-
% Coverage-und-Ledger-Vergleich-6540e2.md; Basis-Urteilsblätter =
% runs/w2-goldv2-eval-ledger-8189ea.json + …-ledger-e30646.json +
% …-f-6540e2.json + …-f-12d046.json (Kopien thesis-evidence/w2/eval/);
% Stufenanalyse = runs/w2-z1-stufenanalyse-8189ea.json (A-Urteile der
% 5 geänderten IDs unter denselben Auslegungen nachgezogen, s.
% konsolidierung.json §aBAbgleich); Signale = runs/ledger/
% 20260905_164650_8189ea/step-01d-unused-unit-ledger-compare/
% output.json; Design-Gruppe = 37 Gold-IDs mit sämtlichen Quell-Units
% ab AU-0117. Läufe: Ledger 20260905_164650_8189ea (Interview) /
% 20260905_164206_e30646 (meeting-2); F 20260906_143052_6540e2 /
% 20260906_141518_12d046. Meeting-2-Urteile unverändert
% (Konsistenzprüfung 9 IDs / 4 Muster, 0 Änderungen).
% ------------------------------------------------------------

\section{Ergebnisqualität der Ledger-Kette}
\label{sec:eval-ledgerqualitaet}

Dieser Abschnitt untersucht den maschinellen Claim-Bestand der
Ledger-Kette vor dem Human-Gate (\emph{Pre-Gate}) gegen den in
Abschnitt~\ref{sec:eval-design} beschriebenen Referenzbestand (Gold~v2: 109 bzw.\ 31
Referenzaussagen). Zunächst werden erhaltene Inhalte und Lücken
bestimmt (Abschnitt~\ref{sec:eval-kernzahlen}). Danach folgen die
Zuordnung zu Prüfhinweisen und die Lokalisierung von Verlusten in der
Kette (Abschnitte~\ref{sec:eval-luecken} und~\ref{sec:eval-stufen}).
Abschließend wird die freie Vergleichskonfiguration F eingeordnet
(Abschnitt~\ref{sec:eval-armf}). Die Abdeckungswerte der Hauptläufe
entsprechen dem konsolidierten Stand vom 08.09.; die ergänzende
Ledger-Auswertung und die erneute Stützungsprüfung vom 15.09. sind
jeweils gekennzeichnet (Abschnitt~\ref{sec:eval-design}). Sämtliche
Bewertungen betreffen unveränderte historische Ausgaben. Ein externer Akzeptanzmaßstab für eine ausreichende Abdeckung wurde
nicht festgelegt; die Zahlen beschreiben den erreichten Stand.

\subsection{Zwei Fälle, zwei Aufgabenprofile}
\label{sec:eval-kernzahlen}

Tabelle~\ref{t:eval-kernzahlen-v2} stellt beide Quellenfälle
nebeneinander. Der \emph{Treue-Fall} (meeting-2) ist ein kompaktes,
klar strukturiertes Meeting; der \emph{Stressfall}
(Interview-Einrichtung) ist ein zusammengeführter Interview- und
Ausarbeitungsverlauf über mehrere Gesprächsphasen, dessen
Referenzbestand zu gut einem Drittel aus späten
UI-Detailfestlegungen besteht (37 von 109 Aussagen mit sämtlichen
Quell-Units im Design-Abschnitt).

% B-59, 15.09.2026: datierte Ledger-Nachauswertung, Stützung und Reichweite.
% Beleg: thesis-evidence/w2/ledger-repetitions-20260915/ (Originale unverändert).
% M19 (12.09.2026): Leserführung und präzise Ergebnisdarstellung;
% Messwerte und historische Nachweisgrenzen bleiben erhalten.
% ============================================================
% Tabelle 7.3-1: Kernzahlen beider Fälle (Gold v2) — v02 (08.09.)
% In Overleaf ablegen als: tables/Chap7-v02/chap7.3table1.tex
% Aufruf im Fließtext: % B-59, 15.09.2026: datierte Ledger-Nachauswertung, Stützung und Reichweite.
% Beleg: thesis-evidence/w2/ledger-repetitions-20260915/ (Originale unverändert).
% M19 (12.09.2026): Leserführung und präzise Ergebnisdarstellung;
% Messwerte und historische Nachweisgrenzen bleiben erhalten.
% ============================================================
% Tabelle 7.3-1: Kernzahlen beider Fälle (Gold v2) — v02 (08.09.)
% In Overleaf ablegen als: tables/Chap7-v02/chap7.3table1.tex
% Aufruf im Fließtext: % B-59, 15.09.2026: datierte Ledger-Nachauswertung, Stützung und Reichweite.
% Beleg: thesis-evidence/w2/ledger-repetitions-20260915/ (Originale unverändert).
% M19 (12.09.2026): Leserführung und präzise Ergebnisdarstellung;
% Messwerte und historische Nachweisgrenzen bleiben erhalten.
% ============================================================
% Tabelle 7.3-1: Kernzahlen beider Fälle (Gold v2) — v02 (08.09.)
% In Overleaf ablegen als: tables/Chap7-v02/chap7.3table1.tex
% Aufruf im Fließtext: \input{tables/Chap7-v02/chap7.3table1}
% Label: t:eval-kernzahlen-v2
% Stil: chap6.1-Konvention.
% HERKUNFT: konsolidierter Auswertungsstand 08.09. =
% runs/w2-nachreview-konsolidierung.json (Nachreview-Runde: 5
% geänderte Ledger-Urteile 025/062/086/105/107); Basis =
% runs/w2-goldv2-eval-ledger-8189ea.json (Interview) +
% runs/w2-goldv2-eval-ledger-e30646.json (meeting-2, unverändert),
% Kopien thesis-evidence/w2/eval/; Nenner = Gold v2 (109/31, Hash ②).
% ============================================================

\begin{table}[htbp]
  \centering
  \small
  \begin{tabular}{>{\raggedright\arraybackslash}p{0.40\textwidth}rr}
    \hline
    & \textbf{Treue-Fall} & \textbf{Stressfall} \\
    & (meeting-2) & (Interview) \\
    \hline
    Quell-Units / Referenzaussagen (Gold v2) & 36 / 31 & 136 / 109 \\
    Gelieferte kanonische Claims & 20 & 47 \\
    Urteile full / partial / none & 26 / 2 / 3 & 54 / 29 / 26 \\
    \hline
    Full Coverage (strict) & 0{,}839 & 0{,}495 \\
    Touched Coverage & 0{,}903 & 0{,}761 \\
    Inhaltliche Korrektheit der Claims & 1{,}0 & 1{,}0 \\
    Semantische Stützung & 1{,}0 & 0{,}957 \\
    \hline
    Lücken gesamt (partial + none) & 5 & 55 \\
    \quad davon mit zuordenbarem Prüfhinweis & 0 & 37 \\
    \quad davon still & 5 & 18 \\
    Silent-Miss-Rate & 1{,}0 & 0{,}327 \\
    Adressierbarkeit (Prüfpotenzial) & 0{,}839 & 0{,}835 \\
    Prüfhinweise gesamt / lückenrelevant & 3 / 0 & 48 / 18 \\
    Signalbezogene Präzision & 0{,}0 & 0{,}375 \\
    \hline
  \end{tabular}
  \caption[Ergebnisqualität der Ledger-Kette vor dem Human-Gate]{Ergebnisqualität der Ledger-Kette (Pre-Gate) gegen den
    Referenzbestand Gold~v2, je Fall ein Lauf gemäß
    dokumentierter Auswahlregel;
    Stressfall-Abdeckung konsolidiert am 08.09.; Stützung am
    15.09. nachgeprüft (45/47, \hyperref[anh:ledger-wiederholungen]{Anhang~D.11}).
    Bewertungsrevisionen betreffen dieselben Ausgaben, keine
    Systemänderung. „Inhaltliche
    Korrektheit`` folgt dem tatsächlichen Urteilsverfahren:
    quellengetragene Inhalte ohne Gold-Gegenstück zählen nicht als
    Fehler (gesonderte \emph{gold\_escapes}-Zählung nur soweit
    erhoben, Anhang~\ref{anh:evaluation}, \hyperref[anh:referenzpruefung]{Abschnitt~D.3}; keine Gold-Precision im engeren Sinn). Adressierbarkeit
    bezeichnet Prüfpotenzial, keine erreichte Abdeckung; jede Lücke
    und jeder Hinweis zählt je Kennzahl nur einmal.}
  \label{t:eval-kernzahlen-v2}
\end{table}

% Label: t:eval-kernzahlen-v2
% Stil: chap6.1-Konvention.
% HERKUNFT: konsolidierter Auswertungsstand 08.09. =
% runs/w2-nachreview-konsolidierung.json (Nachreview-Runde: 5
% geänderte Ledger-Urteile 025/062/086/105/107); Basis =
% runs/w2-goldv2-eval-ledger-8189ea.json (Interview) +
% runs/w2-goldv2-eval-ledger-e30646.json (meeting-2, unverändert),
% Kopien thesis-evidence/w2/eval/; Nenner = Gold v2 (109/31, Hash ②).
% ============================================================

\begin{table}[htbp]
  \centering
  \small
  \begin{tabular}{>{\raggedright\arraybackslash}p{0.40\textwidth}rr}
    \hline
    & \textbf{Treue-Fall} & \textbf{Stressfall} \\
    & (meeting-2) & (Interview) \\
    \hline
    Quell-Units / Referenzaussagen (Gold v2) & 36 / 31 & 136 / 109 \\
    Gelieferte kanonische Claims & 20 & 47 \\
    Urteile full / partial / none & 26 / 2 / 3 & 54 / 29 / 26 \\
    \hline
    Full Coverage (strict) & 0{,}839 & 0{,}495 \\
    Touched Coverage & 0{,}903 & 0{,}761 \\
    Inhaltliche Korrektheit der Claims & 1{,}0 & 1{,}0 \\
    Semantische Stützung & 1{,}0 & 0{,}957 \\
    \hline
    Lücken gesamt (partial + none) & 5 & 55 \\
    \quad davon mit zuordenbarem Prüfhinweis & 0 & 37 \\
    \quad davon still & 5 & 18 \\
    Silent-Miss-Rate & 1{,}0 & 0{,}327 \\
    Adressierbarkeit (Prüfpotenzial) & 0{,}839 & 0{,}835 \\
    Prüfhinweise gesamt / lückenrelevant & 3 / 0 & 48 / 18 \\
    Signalbezogene Präzision & 0{,}0 & 0{,}375 \\
    \hline
  \end{tabular}
  \caption[Ergebnisqualität der Ledger-Kette vor dem Human-Gate]{Ergebnisqualität der Ledger-Kette (Pre-Gate) gegen den
    Referenzbestand Gold~v2, je Fall ein Lauf gemäß
    dokumentierter Auswahlregel;
    Stressfall-Abdeckung konsolidiert am 08.09.; Stützung am
    15.09. nachgeprüft (45/47, \hyperref[anh:ledger-wiederholungen]{Anhang~D.11}).
    Bewertungsrevisionen betreffen dieselben Ausgaben, keine
    Systemänderung. „Inhaltliche
    Korrektheit`` folgt dem tatsächlichen Urteilsverfahren:
    quellengetragene Inhalte ohne Gold-Gegenstück zählen nicht als
    Fehler (gesonderte \emph{gold\_escapes}-Zählung nur soweit
    erhoben, Anhang~\ref{anh:evaluation}, \hyperref[anh:referenzpruefung]{Abschnitt~D.3}; keine Gold-Precision im engeren Sinn). Adressierbarkeit
    bezeichnet Prüfpotenzial, keine erreichte Abdeckung; jede Lücke
    und jeder Hinweis zählt je Kennzahl nur einmal.}
  \label{t:eval-kernzahlen-v2}
\end{table}

% Label: t:eval-kernzahlen-v2
% Stil: chap6.1-Konvention.
% HERKUNFT: konsolidierter Auswertungsstand 08.09. =
% runs/w2-nachreview-konsolidierung.json (Nachreview-Runde: 5
% geänderte Ledger-Urteile 025/062/086/105/107); Basis =
% runs/w2-goldv2-eval-ledger-8189ea.json (Interview) +
% runs/w2-goldv2-eval-ledger-e30646.json (meeting-2, unverändert),
% Kopien thesis-evidence/w2/eval/; Nenner = Gold v2 (109/31, Hash ②).
% ============================================================

\begin{table}[htbp]
  \centering
  \small
  \begin{tabular}{>{\raggedright\arraybackslash}p{0.40\textwidth}rr}
    \hline
    & \textbf{Treue-Fall} & \textbf{Stressfall} \\
    & (meeting-2) & (Interview) \\
    \hline
    Quell-Units / Referenzaussagen (Gold v2) & 36 / 31 & 136 / 109 \\
    Gelieferte kanonische Claims & 20 & 47 \\
    Urteile full / partial / none & 26 / 2 / 3 & 54 / 29 / 26 \\
    \hline
    Full Coverage (strict) & 0{,}839 & 0{,}495 \\
    Touched Coverage & 0{,}903 & 0{,}761 \\
    Inhaltliche Korrektheit der Claims & 1{,}0 & 1{,}0 \\
    Semantische Stützung & 1{,}0 & 0{,}957 \\
    \hline
    Lücken gesamt (partial + none) & 5 & 55 \\
    \quad davon mit zuordenbarem Prüfhinweis & 0 & 37 \\
    \quad davon still & 5 & 18 \\
    Silent-Miss-Rate & 1{,}0 & 0{,}327 \\
    Adressierbarkeit (Prüfpotenzial) & 0{,}839 & 0{,}835 \\
    Prüfhinweise gesamt / lückenrelevant & 3 / 0 & 48 / 18 \\
    Signalbezogene Präzision & 0{,}0 & 0{,}375 \\
    \hline
  \end{tabular}
  \caption[Ergebnisqualität der Ledger-Kette vor dem Human-Gate]{Ergebnisqualität der Ledger-Kette (Pre-Gate) gegen den
    Referenzbestand Gold~v2, je Fall ein Lauf gemäß
    dokumentierter Auswahlregel;
    Stressfall-Abdeckung konsolidiert am 08.09.; Stützung am
    15.09. nachgeprüft (45/47, \hyperref[anh:ledger-wiederholungen]{Anhang~D.11}).
    Bewertungsrevisionen betreffen dieselben Ausgaben, keine
    Systemänderung. „Inhaltliche
    Korrektheit`` folgt dem tatsächlichen Urteilsverfahren:
    quellengetragene Inhalte ohne Gold-Gegenstück zählen nicht als
    Fehler (gesonderte \emph{gold\_escapes}-Zählung nur soweit
    erhoben, Anhang~\ref{anh:evaluation}, \hyperref[anh:referenzpruefung]{Abschnitt~D.3}; keine Gold-Precision im engeren Sinn). Adressierbarkeit
    bezeichnet Prüfpotenzial, keine erreichte Abdeckung; jede Lücke
    und jeder Hinweis zählt je Kennzahl nur einmal.}
  \label{t:eval-kernzahlen-v2}
\end{table}


\emph{Quellenbindung und Inhaltsurteil.} In den beiden
Hauptausgaben wurde bei keinem der 67 Claims eine Quellenverfälschung
festgestellt. Vollständig durch die angegebenen Units gestützt sind
nach dem jeweiligen Prüfstand 20/20 Claims im Treue-Fall und 45/47
im Stressfall. Dessen frühere Bewertung von 47/47 wurde am 15.09.
berichtigt: Bei zwei Claims liegt notwendiger Kontext außerhalb
der referenzierten Units. Beispielsweise nennt ein Claim Animationen,
verweist aber nur auf eine Antwort, deren Gegenstand erst aus der
vorhergehenden Frage hervorgeht. Der Inhalt bleibt im Transkript
auffindbar; die angegebene Textstelle trägt ihn allein nicht
vollständig (\hyperref[anh:ledger-wiederholungen]{Anhang~D.11}).
Quellengetragene Inhalte ohne Gold-Gegenstück (etwa Kontext- und
Prozessaussagen) zählen verfahrensgemäß nicht als Fehler. Sie werden, soweit
erhoben, separat als \emph{gold\_escapes} geführt. Die Korrektheitskennzahl
misst damit Verfälschungsfreiheit gegenüber der Quelle (Definition,
Erhebungsstand und dokumentierte Protokoll-Abweichung:
Abschnitt~\ref{sec:eval-design}). Die
technische Auflösbarkeit ist deterministisch geprüft: Sämtliche abgegebenen
Einzelreferenzen lösen auf: beim Ledger 33/33 im Treue-Fall und
87/87 im Stressfall, bei F entsprechend 44/44 und 208/208.
In keinem der vier Läufe ist ein Claim referenzlos.
Eine nach der Gold-Revision geprüfte Zählung quellengetragener
Inhalte außerhalb des Referenzbestands liegt nur für den
F-Treue-Fall vor (drei, alle gestützt); für die übrigen Läufe
wird keine konsolidierte Zählung berichtet --- ein historischer
Altzähler des F-Stressfall-Blatts auf dem früheren Referenzstand
wird nicht übernommen (Anhang~\ref{anh:evaluation}).

\emph{Abdeckung.} Der Stressfall-Bestand deckt knapp die Hälfte der
Referenzaussagen vollständig (0{,}495 strict), der Treue-Fall
0{,}839. Dabei ist die ungewichtete Aussageabdeckung wörtlich zu
nehmen: Jede Referenzaussage zählt gleich viel, und eine als
unvollständig oder nicht abgedeckt gewertete Aussage bedeutet oft
einen fehlenden Teilaspekt \emph{innerhalb} einer erkannten
Funktion --- auch ein „none`` belegt nicht, dass ein komplettes
Modul fehlt (Beispiel: das Kalendermodul ist erfasst, das
Hinzufügen von Tageseinträgen nicht). Im Design-Abschnitt des
Stressfalls sind 26 von 37 Referenzaussagen unvollständig
abgedeckt, im übrigen Referenzbestand 29 von 72 --- die späten
UI-Festlegungen bilden einen überproportionalen Schwerpunkt der
beobachteten Lücken, erklären die Differenz zum Treue-Fall aber
nicht allein.

\emph{Prüfhinweise.} Im Stressfall ließ sich 37 der 55 Lücken
mindestens ein Hinweis zuordnen. Zusammen mit den 54 vollständig
abgedeckten Aussagen ergibt das eine Adressierbarkeit von
$(54+37)/109 = 0{,}835$. Diese Kennzahl beschreibt das
\emph{Prüfpotenzial}; sie misst weder zusätzliche Abdeckung noch
menschliche Behebung. Im Treue-Fall blieben alle fünf Lücken ohne
Hinweis, weil sie sämtlich in bereits verwendeten Units liegen.

\emph{Ergänzende Wiederholungsprüfung.} Zwei weitere archivierte
Stressfall-Ausgaben erreichen 57 beziehungsweise 54 Volltreffer.
Zusammen mit dem Hauptlauf ergibt sich eine beobachtete Spanne von
49{,}5--52{,}3\,\% Full Coverage bei einem Median von 49{,}5\,\%.
Die ähnlichen Gesamtwerte verdecken Unterschiede im erhaltenen
Inhalt: Nur 43 Referenzaussagen sind in allen drei Ausgaben
vollständig enthalten; bei 25 wechselt der Status zwischen
vollständig und nicht vollständig. Beide Zusatzläufe
enthalten ausdrückliche Verbote für Nutzer, Daten zu löschen oder
Accounts anzulegen. Diese fehlen im Hauptlauf. Dessen Entscheidung zur
Auslagerung vertiefter Datenschutzplanung fehlt dagegen in den
Zusatzläufen.
Alle Quellenreferenzen sind technisch gültig; die Stützungsprüfung
ergibt für die Zusatzläufe 44/47 und 38/41 vollständig gestützte
Claims. Jeweils ein Claim schwächt zudem den Planungsstand der Quelle
ab. Die drei Ausgaben beschreiben diesen Quellenfall unter den
protokollierten Bedingungen. Hinweise, menschliche Behebung und
allgemeine Zuverlässigkeit wurden damit nicht zusätzlich untersucht
(\hyperref[anh:ledger-wiederholungen]{Anhang~D.11}).

\subsection{Lücken und Prüfhinweise}
\label{sec:eval-luecken}

\textbf{a) Was fehlt inhaltlich?} Die 55 unvollständig
abgedeckten Referenzaussagen des Stressfalls wurden retrospektiv
nach dem jeweils \emph{fehlenden Bestandteil} typisiert
(Tabelle~\ref{t:eval-lueckentypen}; qualitative Inhaltsbeschreibung,
keine validierte Wichtigkeitsskala): 18-mal fehlt eine Funktion,
Berechtigung oder ein fachlicher Ablauf, 18-mal Darstellung oder
konkrete Bediengestaltung, 11-mal ein Zweck- oder Qualitätsziel,
8-mal eine offene Frage oder technische Konkretisierung. Drei
belegte Beispiele zeigen die Spannweite: Bei der Profilanlage
fehlt nur der Bedienweg (Plus-Symbol unten samt Dialog --- die
Funktion selbst ist erhalten); beim Berechtigungsmodell fehlt die
kurze, fachlich gewichtige Verbotsklausel „User dürfen nichts
löschen``; beim Kommunikationswissen fehlt die
Qualitätszusage des jederzeit schnellen Nachschlagens. Eine kurze
Formulierung kann fachlich wesentlich sein --- die Typen sagen
nicht, dass Darstellungs- oder Zweck-Lücken unwichtig wären.

\textbf{b) Welche Prüfansätze entstehen --- und welche erreichen
die Vorlage?} Die deterministische Bilanz des Stresslaufs erfasste
75 der 136 Quell-Units mit Bezug zu mindestens einem Kandidaten;
61 blieben ohne Kandidatenbezug. Alle 61 wurden modellgestützt triagiert,
davon 48 als potenziell relevant eingeordnet und gegen den
Kandidatenbestand geprüft. Daraus entstanden 24 Vorschläge zum
Anhängen als Beleg, 21 Einstufungen als mittelbar abgedeckt, ein
Hinweis auf einen fehlenden Claim und zwei menschlich zu klärende
Fälle. Die gespeicherten Begründungen machen diese Einordnungen
nachlesbar. Die 24 Belegvorschläge waren am Messpunkt vor dem
Human-Gate keine ausgeführten Ergänzungen. Die 48 Hinweise sind
damit Verarbeitungsergebnisse, nicht 48 erkannte Inhaltslücken
(Abschnitt~\ref{sec:eval-kontrollen} nutzt ihre Herkunft für die Kontrollbetrachtung).
Als einem Hinweis zugeordnet gilt eine Lücke nach einer festen
Regel: wenn die Hinweis-Unit in den Quell-Units der
Referenzaussage vorkommt (Unit-Schnittmenge); eine einzige,
dokumentierte semantische Zusatz-Zuordnung ergänzt diese Regel.

18 der 48 Hinweise sind
für mindestens eine tatsächliche Lücke relevant (signalbezogene
Präzision 0{,}375); 37 der 55 Lücken tragen mindestens eine
Zuordnung (ohne die semantische Ergänzung 36; Silent-Miss dann
19/55 statt 18/55).

In die rekonstruierte Adjudikations-Queue (45
Positionen: 18 Claim-Reviews, 27 Unit-Hinweise) gelangen davon
\emph{29 Lücken} --- acht weitere verlieren ihre Vorlage, weil ihr
einziger Hinweis das Prüfergebnis \texttt{already\_covered\_indirectly}
trägt, also „indirekt bereits abgedeckt``, und der Queue-Aufbau
solche Hinweise verwirft. Die Einzelprüfung zeigt: Die Einstufungen
treffen das Kernthema, übergehen aber die fehlende Qualifikation.
Diese Lücken entfallen damit aus dem Unit-Hinweiskanal; andere
Zugänge zur Prüfung sind dadurch nicht ausgeschlossen. Nur ein
einziges Queue-Item beanstandet ein Fehlen \emph{explizit} samt
Formulierungsvorschlag (\texttt{missing\_claim} zur nicht
verwendeten Unit AU-0124, Profil-Detailansicht --- es trägt zwei
direkte Gold-Bezüge und die eine semantische Zuordnung); die
übrigen Hinweise schlagen beruhigendere Einstufungen vor.
Nur ein Teil der Hinweise verweist somit auf tatsächliche Lücken;
ihre Einstufung benennt das Fehlen meist nicht ausdrücklich.

Ein zweiter Zugang existiert über die 18
Claim-Reviews: Deren gespeicherte Zitate enthalten in Einzelfällen
genau den im Aussagentext fehlenden Inhalt, etwa die Profilanlage
per Plus-Symbol und Dialog. Dieser andere Zugang erhöht die
Unit-Signal-Zählung nicht; eine tatsächliche Nutzung der Belege zur
inhaltlichen Ergänzung ist nicht nachgewiesen.

Beispiel für einen \emph{stillen Detailverlust in einer
verwendeten Unit}: AU-0102 (einheitliche Flutter-/Dart-Versionen)
wurde von einem Claim referenziert, der die
Versionsvereinheitlichung trägt --- die konkreten Versionsnummern
3.13.9/3.1.5 ließ er fallen; da die Unit als „verwendet`` galt,
entstand kein Hinweis. Alle 18 stillen Lücken des Stressfalls
(und alle fünf des Treue-Falls) gehören zu dieser Klasse.

\textbf{c) Was wäre unter idealisierter Ergänzung möglich?} Als
ausdrücklich hypothetische Rechnung: Würden alle 29 in der Queue
mit einem Unit-Hinweis verknüpften Lücken vollständig geschlossen
--- was voraussetzt, dass der gesamte fehlende Gehalt erkannt
wird, der erforderliche Quellenkontext zugänglich ist, eine
zulässige Bearbeitung den Aussagentext vollständig ergänzt und
alle bisherigen Volltreffer erhalten bleiben ---, ergäbe sich eine
bedingte Full Coverage von $(54+29)/109 = 0{,}761$. Das ist keine
gemessene Post-HITL-Abdeckung, keine Wahrscheinlichkeit und keine
bewiesene Obergrenze der Behebbarkeit: Mehrere zugeordnete Units
tragen den fehlenden Gehalt nicht allein (AU-0024 trägt die
Nutzenaussagen zweier Lücken nicht; AU-0118 trägt bei der
appweiten Usability-Aussage zunächst nur den
Eingabefeld-Kontext). Über die 26 nicht angesetzten Lücken trifft
die Rechnung keine Aussage zu weiteren Zugängen; für einzelne
dieser Restfälle sind Claim-Review-Kontexte belegt, eine Vorlage
aller 26 ist nicht gezeigt. Der Zahlenwert stimmt
zufällig mit der Touched Coverage überein; die beiden 29er-Mengen
sind nicht identisch (das Szenario zählt 13 teilweise und 16
nicht abgedeckte Lücken).

\subsection{Begrenzte Stufenanalyse: Wo entsteht der Verlust?}
\label{sec:eval-stufen}

Für den Stresslauf wurde derselbe Referenzmaßstab an zwei
Beständen desselben Laufs angelegt: an den 54 Extraktions-Kandidaten
(Bestand~A, \texttt{step-01}) und am kanonischen Endbestand
(Bestand~B, 47 Claims). Die A-Urteile der fünf in der
Nachreview-Runde geänderten Referenzaussagen wurden unter
denselben Auslegungen nachgezogen (alle fünf folgen ihrer
B-Klasse; Belege im Konsolidierungs-Artefakt).
Tabelle~\ref{t:eval-stufen} zeigt den Übergang; jede
Referenzaussage belegt genau eine Zelle.

% P10-4b, 14.09.2026: Begriffs-/Beleganschlüsse und abgegrenzte Prüfaussagen präzisiert.
% ============================================================
% Tabelle 7.3-2: 3x3-Übergangstabelle der Stufenanalyse — v02 (08.09.)
% In Overleaf ablegen als: tables/eval-stufen-uebergang.tex
% Aufruf im Fließtext: \input{\tabledir/eval-stufen-uebergang.tex}
% Label: t:eval-stufen
% Stil: chap6.1-Konvention.
% HERKUNFT: runs/w2-z1-stufenanalyse-8189ea.json (alle 109 A-/B-
% Urteile mit Beleg) + runs/w2-nachreview-konsolidierung.json
% §aBAbgleich (A-Urteile der 5 geänderten IDs 025/062/086/105/107
% unter denselben Auslegungen nachgezogen — alle folgen ihrer
% B-Klasse; Klassenwechsel bleiben exakt G-IE-021/023).
% Lauf: 20260905_164650_8189ea.
% ============================================================

\begin{table}[htbp]
  \centering
  \small
  \begin{tabular}{lrrrr}
    \hline
    & \multicolumn{3}{c}{\textbf{Endbestand B (47 Claims)}} & \\
    \textbf{Kandidaten A (54 Claims)} & full & partial & none & $\Sigma$ A \\
    \hline
    full    & 54 & 1 & 1 & 56 \\
    partial & 0 & 28 & 0 & 28 \\
    none    & 0 & 0 & 25 & 25 \\
    \hline
    $\Sigma$ B & 54 & 29 & 26 & 109 \\
    \hline
  \end{tabular}
  \caption[Abdeckung vor und nach der Kanonisierung im Stresslauf]{Übergang der Abdeckungsklassen aller 109 Referenzaussagen
    zwischen Extraktions-Kandidaten (A) und kanonischem Endbestand
    (B) desselben Stresslaufs, im konsolidierten Auswertungsstand
    (A-Urteile der fünf neu entschiedenen Aussagen unter denselben
    Auslegungen nachgezogen). Nur zwei Aussagen wechseln die Klasse
    (beide Verluste, beide aus derselben Zusammenführung); die
    kleinere Claim-Zahl selbst ist kein Verlustnachweis.}
  \label{t:eval-stufen}
\end{table}


Im untersuchten Übergang behielten 107 der 109 Referenzaussagen
ihre Abdeckungsklasse; Verbesserungen traten nicht auf. „Gleiche
Klasse`` ist dabei die präzise Aussage --- ein
partial-zu-partial-Übergang kann Veränderungen innerhalb der
Klasse enthalten. Zwei zuvor vollständig enthaltene Aussagen
wurden nur noch teilweise beziehungsweise nicht mehr abgedeckt
--- beide betrafen dieselbe kanonische Zusammenführung: Der
Kandidat DEC-019 enthielt wörtlich „User können Inhalte
hinzufügen, aber nichts löschen``; der aus DEC-019 und DEC-020
zusammengeführte kanonische Claim übernahm Rollen, Admin-Rechte
und Registrierungsverbot, ließ diese Klausel jedoch fallen, und
kein anderer Claim des Endbestands trägt sie. Der referenzielle
Übergangscheck meldete dabei keinen Fehler, denn der Claim
referenziert beide Kandidaten korrekt --- für diese beiden
Inhalte war eine gültige Kandidatenverknüpfung somit kein
ausreichender Nachweis semantischer Erhaltung. Die Mengenänderung
von 54 auf 47 Claims ist dagegen kein Verlustnachweis: Sie geht
vollständig auf sieben Zweiergruppen zurück, von denen sechs ohne
Klassenverlust der geprüften Referenzinhalte blieben (was
semantische Verlustfreiheit im Detail nicht beweist).

Damit ist die Abdeckungsgrenze der Kette lokalisiert: 53 der 55
Lücken des Endbestands bestanden bereits im Extraktions-Bestand~A
und wurden im untersuchten weiteren Verlauf nicht vollständig
geschlossen; zwei zusätzliche Abdeckungsverluste entstanden
zwischen A und B. Die gespeicherten Zwischenstände machen genau
diese Lokalisierung möglich --- das ist die belegte Eigenschaft;
verhindert hat der referenzielle Übergangscheck den semantischen
Verlust nicht. Der Befund beschreibt den Gesamtübergang dieses
einen Laufs (Kanonisierung einschließlich Übergangs-Check); eine
Zuschreibung an einen einzelnen Modellaufruf tragen die Artefakte
nicht.

\subsection{Ergänzende Einordnung: die freie Vergleichsbedingung F}
\label{sec:eval-armf}

Als Kontrastfolie verarbeitete eine Einzel-Agent-Konfiguration~F
dieselben Fälle: gleiche Quelle, ausdrücklich beauftragte
konsolidierte, atomare, kanonische Aussagen mit erhaltenen
Qualifikationen und vollständiger Quellenbilanz
(\texttt{used}/\texttt{non\_relevant}/\texttt{unresolved} je
Unit), freie Arbeitsweise mit genutzter Vorprüfung über mehrere
Modellrunden --- damit geht F über einen einzelnen Roh-Extraktionsaufruf hinaus.
Anders als der frühe Erzeugungsansatz aus
Abschnitt~\ref{sec:evidenzgebundene-wissensgewinnung} verlangt F
bereits explizite Quellenreferenzen und eine Bilanz. Der Vergleich
prüft somit die bezeichneten späteren Konfigurationen, nicht
unverändert den damaligen Ausgangsversuch.

F erreichte im Stressfall die höhere direkte Abdeckung
(0{,}633 strict / 0{,}844 touched) bei rund 78\,\% der Eingabe-
und 48\,\% der Ausgabe-Tokens des Ledger-Laufs. Die inhaltliche
Korrektheit beträgt 0{,}991; es gilt dieselbe Kennzahl wie in
Tabelle~\ref{t:eval-kernzahlen-v2}. Ihre claimbezogenen Quoten sind
wegen der unterschiedlichen Granularität jedoch nur eingeschränkt
vergleichbar: 47 gebündelte Claims stehen 117 feingranularen Aussagen
gegenüber. Aus 1{,}0 gegenüber 0{,}991 folgt deshalb kein belastbarer
genereller Qualitätsvorsprung. Die 117 Aussagen sind eine
Bestandsgröße, keine Qualitätskennzahl.

Der gepaarte Vergleich der beiden Stressfall-Hauptläufe über dieselben
109 Referenzaussagen ergibt: 49 sind bei beiden vollständig erhalten,
20 nur bei F, fünf nur beim Ledger, 35 bei keinem (vollständige
ID-Mengen im Anhang). Die 20 nur von F gedeckten Fälle sind
\emph{nicht} überwiegend UI-Details: Nach dem beim Ledger
fehlenden Bestandteil verteilen sie sich auf acht funktionale
beziehungsweise rollenbezogene, fünf Darstellungs-, vier
Zweck-/Qualitäts- und drei offene/technische Aspekte ---
beispielsweise trägt F das ausdrückliche Verbot, dass normale
User hinzugefügte Daten löschen (ein anderer Gegenstand als das
bei beiden fehlende Verbot der Account-Anlage), und die
leicht zugänglichen Datenschutzeinstellungen. Umgekehrt erhält
der Ledger unter anderem die Klärung der Dokumentationsübernahme
mit der Einrichtungsleitung und die Entscheidung, die vertiefte
Datenschutzkonzeption als Folgeprojekt auszulagern --- beides
fehlt im F-Hauptlauf. Diese Beispiele kennzeichnen Unterschiede
dieses Laufpaars. In beiden zusätzlichen Ledger-Ausgaben ist das
Löschverbot enthalten, die Auslagerung der Datenschutzkonzeption
dagegen nicht (Abschnitt~\ref{sec:eval-kernzahlen}). Daraus lassen
sich keine beständigen inhaltlichen Stärken der Konfigurationen
ableiten; für F liegt keine entsprechende Wiederholungsprüfung vor.
Die Lückenprofile der Hauptläufe ähneln sich strukturell
(Tabelle~\ref{t:eval-lueckentypen}): Beide Konfigurationen
verlieren überwiegend Qualifikationen, Gestaltungsmerkmale und
Zweckaussagen innerhalb erkannter Themen.

% ============================================================
% Tabelle 7.3-3: Lückentypen beider Konfigurationen — v01 (08.09.)
% In Overleaf ablegen als: tables/eval-lueckentypen.tex
% Aufruf im Fließtext: \input{\tabledir/eval-lueckentypen.tex}
% Label: t:eval-lueckentypen
% Stil: chap6.1-Konvention.
% HERKUNFT: runs/w2-nachreview-konsolidierung.json §lueckentypen
% (konsolidierter Stand; Typisierung je ID aus den beiden Reviews
% übernommen, drei neue Lücken dokumentiert nachtypisiert:
% Ledger-062=F, 107=O beidseitig, F-089=U). Typ beschreibt den
% FEHLENDEN Bestandteil der jeweiligen Ausgabe — dieselbe Gold-ID
% kann je Konfiguration verschieden typisiert sein.
% ============================================================

\begin{table}[htbp]
  \centering
  \small
  \begin{tabular}{>{\raggedright\arraybackslash}p{0.52\textwidth}rr}
    \hline
    \textbf{Fehlender Bestandteil} & \textbf{Ledger} & \textbf{F} \\
    & (55 Lücken) & (40 Lücken) \\
    \hline
    Funktion / Berechtigung / fachlicher Ablauf & 18 & 6 \\
    Darstellung / konkrete Bediengestaltung & 18 & 18 \\
    Zweck / Qualitätsziel & 11 & 8 \\
    Offene Frage / technische Konkretisierung / Interpretationsfall & 8 & 8 \\
    \hline
    Summe & 55 & 40 \\
    \hline
  \end{tabular}
  \caption[Typisierung der Abdeckungslücken im Stressfall]{Retrospektive Typisierung der unvollständig abgedeckten
    Referenzaussagen des Stressfalls nach dem jeweils fehlenden
    Bestandteil (konsolidierter Auswertungsstand). Der Typ
    beschreibt, \emph{was} in der jeweiligen Ausgabe fehlt --- er
    ist keine Wichtigkeits- oder Schweregradskala, und dieselbe
    Referenzaussage kann je Konfiguration unterschiedlich typisiert
    sein, wenn die Ausgaben unterschiedliche Teile erhalten.}
  \label{t:eval-lueckentypen}
\end{table}


\emph{Umgang mit eigenen Lücken.} F meldete über den vorgesehenen \texttt{unresolved}-Kanal keine
einzige Lücke (bei 40 realen Lücken; Bilanz formal vollständig,
101 verwendete und 35 als nicht relevant eingestufte Units). Die
17 mitgelieferten Unsicherheitsmarker kennzeichnen bestehende
Aussagen als unsicher, benennen aber keine fehlenden Inhalte; die
Quellen-IDs eröffnen menschliches Nachlesen, eine systematische
Vorlage der fehlenden Inhalte gibt es nicht. Die
100-Prozent-Silent-Miss-Angabe gilt entsprechend nur für diese
Signalregel und beweist keine dauerhafte Unauffindbarkeit.

Im Ledger-Bestand konnten dagegen 37 der 55 unvollständig erfassten
Referenzaussagen nachträglich mindestens einem Hinweis zugeordnet
werden --- ein Ansatzpunkt für weitere Prüfung, weder eine
explizite Lückenmeldung noch ein Nachweis tatsächlicher
menschlicher Erkennung oder Behebung; die breitere
Ledger-Zuordnung über Quell-Units ist keine bessere
\emph{explizite} Fehlererkennung.

Ein technischer Nebenbefund
begrenzt zudem die Lesart einer vollkommen konsistenten
F-Bilanz: Alle 208 Einzelreferenzen sind gültig, aber drei
Claim--Unit-Kanten widersprechen sich zwischen Aussagen- und
Bilanzseite (Anhang~\ref{anh:evaluation}); ein spiegelbildlicher
Konsistenznachweis für den Ledger wurde nicht geführt.

Direkte Abdeckung und Prüfpotenzial bleiben getrennte Größen; nicht
beauftragte Dispositionen oder Facetten werden F nicht als Fehler
zugerechnet, und die unterschiedliche Architektur (Einzel-Agent
gegen mehrstufige Kette mit mehr Modellaufrufen und
Zwischenprüfungen) bleibt Teil der Konfiguration --- der
Vergleich isoliert keinen einzelnen Mechanismus kausal. Eine
ergänzende Modellreihe (vier Modelle je Konfiguration; formale
Fertigstellung und semantische Qualität getrennt betrachtet)
findet sich im Anhang.
